\chapter{Conclusioni aperte}
\label{cap:conclusioni}

Questo libro non si chiude con una scoperta. Si chiude, piuttosto, con una domanda meglio posta di quanto non fosse all'inizio, e con gli strumenti per affrontarla con maggiore cognizione di causa.

Il percorso compiuto in queste pagine ha attraversato territori diversi: il confronto tra due paesaggi, la Región de Atacama e le Alpi venete, con le sue piste, le sue coste, le sue miniere e il suo deserto interno costellati di segni incisi o dipinti; il meccanismo per cui una mappa della conoscenza può rispecchiare la storia della ricerca più che la realtà del passato, un meccanismo documentato non in astratto ma in due casi reali, uno sulle Alpi svizzere e uno, ancora più vicino, nelle Dolomiti venete stesse; la lunga vicenda climatica, ricostruita fase per fase con le sue date e i suoi margini di incertezza dichiarati, che ha reso e poi mantenuto abitabili le Alpi venete; il comportamento mutevole della fauna che le ha popolate, con le cifre precise di quanto piccolo fosse il territorio di un camoscio o di uno stambecco, e di quanto quella fedeltà al proprio spazio abbia smentito un'intuizione facile sui loro spostamenti; i limiti di un modello fisico, per quanto solido, applicato a un contesto umano che non gli appartiene, il Qhapaq Ñan e i suoi tratti scavati nella roccia; e infine un metodo che, per quanto nato altrove, sul bacino dell'Adige e dell'Isarco, offre un punto di partenza concreto per chi volesse cercare, con più metodo di quanto sia stato fatto finora, le tracce lasciate sulle Alpi venete da chi le ha attraversate per prima.

Il capitolo precedente ha mostrato che questo non è solo un auspicio teorico. Il Cadore, la parte del Veneto geograficamente più vicina ai casi discussi in questo libro, è passato in pochi anni da area descritta come priva di testimonianze preistoriche a territorio con una frequentazione mesolitica documentata in più punti, semplicemente perché qualcuno ha deciso di cercarla con la stessa pazienza dedicata altrove. È la prova più diretta, tra quelle raccolte in queste pagine, che il sospetto da cui questo libro è partito non è infondato.

Se un solo insegnamento merita di restare, al termine di questo percorso, è forse questo: la scarsità di una traccia non è mai, da sola, una prova della sua assenza. Può esserlo, ma prima di concluderlo occorre aver guardato con la stessa cura con cui si è guardato altrove, dove quella traccia è stata trovata in abbondanza. La Región de Atacama insegna cosa succede quando quello sguardo è stato dedicato, con continuità, per generazioni, a ogni tipo di ambiente, dalle piste di quota al deserto più remoto. Le Alpi venete, in alcuni luoghi, come il Cansiglio e il Cadore, hanno già ricevuto quello sguardo, e i risultati non hanno tardato ad arrivare. In altri, come il massiccio del Grappa, aspettano ancora che qualcuno lo dedichi loro con la stessa pazienza.
